% !TEX root = ../../main.tex
\section{Tổng quan về Hexagonal Architecture}
\label{sec:hexagonal}

\subsection{Hexagonal Architecture là gì?}
\label{subsec:hexagonal_la_gi}
Kiến trúc Lục giác (\textit{Hexagonal Architecture}), còn được gọi là mẫu kiến trúc Cổng và Bộ điều hợp (\textit{Ports and Adapters Pattern}), do chuyên gia phần mềm Alistair Cockburn khởi xướng vào năm 2005 \cite{cockburn2005}. Mục tiêu cốt lõi của kiến trúc này là giải quyết triệt để bài toán phụ thuộc chặt chẽ (\textit{tight coupling}) giữa logic nghiệp vụ trung tâm của phần mềm với các công nghệ giao vận bên ngoài như hệ quản trị cơ sở dữ liệu, giao thức mạng, và giao diện người dùng.

Trong các kiến trúc phân tầng truyền thống (như mô hình 3 tầng Presentation - Business Logic - Data Access), mã nguồn nghiệp vụ thường bị phụ thuộc trực tiếp vào cơ sở dữ liệu hoặc khung giao diện. Ngược lại, Hexagonal Architecture thiết lập một ranh giới bảo vệ nghiêm ngặt: logic nghiệp vụ học thuật nằm tại trung tâm và hoàn toàn cô lập, không mang bất kỳ sự ràng buộc nào với các thư viện ngoại vi, hệ quản trị cơ sở dữ liệu hay giao thức mạng.

Các lợi ích nền tảng của Hexagonal Architecture bao gồm:
\begin{itemize}
    \item \textbf{Độc lập công nghệ:} Cho phép thay thế hệ quản trị cơ sở dữ liệu, chuyển đổi nhà cung cấp mô hình trí tuệ nhân tạo (LLM), hoặc đổi mới giao diện người dùng mà không cần sửa đổi bất kỳ dòng mã logic nghiệp vụ nào.
    \item \textbf{Kiểm thử tự động dễ dàng (Testability):} Logic nghiệp vụ có thể được kiểm thử hoàn toàn độc lập với các cổng giả lập (\textit{Mock Ports}) trong bộ nhớ mà không cần kết nối mạng hay khởi động cơ sở dữ liệu thực tế.
    \item \textbf{Khả năng bảo trì và tiến hóa bền vững:} Ranh giới phân định rõ ràng giúp kiến trúc thích ứng linh hoạt trước sự thay đổi công nghệ nhanh chóng của các nền tảng AI.
\end{itemize}

\subsection{Kiến trúc Ports and Adapters}
\label{subsec:ports_and_adapters}
Hệ thống theo Kiến trúc Lục giác được cấu trúc thành ba phân vùng độc lập tuân thủ nghiêm ngặt nguyên lý đảo ngược phụ thuộc (\textit{Dependency Inversion Principle}):
\begin{enumerate}[label=\textbf{(\arabic*)}]
    \item \textbf{Tầng lõi miền nghiệp vụ thuần khiết (\textit{Pure Core Domain}):} Nằm tại trung tâm của hệ thống, chỉ chứa các thực thể nghiệp vụ (\textit{Entities}), quy tắc logic và các thuật toán tính toán học thuật. Vùng này tuân thủ nguyên tắc cô lập tuyệt đối: chỉ sử dụng các cấu trúc dữ liệu bản địa của ngôn ngữ lập trình, hoàn toàn không phụ thuộc vào bất kỳ thư viện giao tiếp mạng, web framework hay hệ quản trị cơ sở dữ liệu nào.
    \item \textbf{Tầng các cổng giao tiếp trừu tượng (\textit{Ports}):} Là các giao diện trừu tượng (\textit{Interfaces}) nằm trong tầng lõi nghiệp vụ. Cổng quy định hợp đồng giao tiếp chuẩn mà các thành phần bên ngoài bắt buộc phải tuân thủ để phục vụ nghiệp vụ (chẳng hạn như: cổng phân giải ấn phẩm học thuật, cổng bóc tách tài liệu PDF, hay cổng lưu trữ dữ liệu). Có hai loại cổng chính:
    \begin{itemize}
        \item \textit{Cổng phía nhập (Inbound/Driving Ports):} Định nghĩa các phương thức cho phép tác nhân bên ngoài kích hoạt quy trình nghiệp vụ của hệ thống.
        \item \textit{Cổng phía xuất (Outbound/Driven Ports):} Định nghĩa các giao diện cho phép lõi nghiệp vụ yêu cầu tài nguyên hoặc lưu trữ từ hạ tầng bên ngoài.
    \end{itemize}
    \item \textbf{Tầng các bộ điều hợp kỹ thuật (\textit{Adapters}):} Nằm ở lớp ngoài cùng của hệ thống, hiện thực hóa các Cổng trừu tượng để kết nối với các công nghệ thực tế (chẳng hạn như bộ điều hợp cơ sở dữ liệu quan hệ, bộ điều hợp dịch vụ tìm kiếm học thuật trực tuyến, hoặc bộ điều hợp mô hình trí tuệ nhân tạo).
\end{enumerate}

Mô hình phân tầng kiến trúc và ranh giới phân tách trừu tượng của hệ thống được mô tả theo dạng các khối năng lực module trong Hình~\ref{fig:hexagonal_arch}.

\begin{figure}[htbp]
\centering
\begin{tikzpicture}[
    scale=0.74, transform shape,
    card_driving/.style={draw=orange!80!black, fill=orange!3, rounded corners=8pt, line width=1pt, inner sep=8pt},
    card_core/.style={draw=blue!85!black, fill=blue!3, rounded corners=8pt, line width=1.1pt, inner sep=8pt},
    card_driven/.style={draw=teal!85!black, fill=teal!3, rounded corners=8pt, line width=1pt, inner sep=8pt},
    compbox/.style={draw=blue!80!black, fill=white, rounded corners=5pt, line width=0.8pt, align=center, inner sep=4pt, font=\linespread{1.0}\selectfont\footnotesize},
    portbox/.style={draw=teal!80!black, fill=teal!10, rounded corners=5pt, line width=0.9pt, align=center, inner sep=4pt, font=\linespread{1.0}\selectfont\footnotesize},
    adaptbox_dr/.style={draw=orange!80!black, fill=white, rounded corners=5pt, line width=0.8pt, align=center, inner sep=4pt, font=\linespread{1.0}\selectfont\footnotesize},
    adaptbox_dn/.style={draw=teal!80!black, fill=white, rounded corners=5pt, line width=0.8pt, align=center, inner sep=4pt, font=\linespread{1.0}\selectfont\footnotesize},
    arrowstyle/.style={-{Stealth[length=2.5mm]}, line width=1.1pt, draw=blue!85!black},
    dasharrow/.style={-{Stealth[length=2.5mm]}, dashed, line width=1.0pt, draw=teal!85!black}
]
    % --------------------------------------------------------------------------
    % Khối 1: Tầng Điều Hợp Điều Khiển (Driving Adapters) - Bên trái
    % --------------------------------------------------------------------------
    \node[card_driving, minimum width=4.5cm, minimum height=8.6cm] (c_driving) at (-6.0, 0) {};
    \node[align=center, font=\linespread{1.0}\selectfont\bfseries\scriptsize, text=orange!90!black] at (-6.0, 3.8) {BỘ ĐIỀU HỢP ĐIỀU KHIỂN\\[2pt](DRIVING ADAPTERS)};
    
    \node[adaptbox_dr, text width=3.6cm, minimum width=3.8cm] (ad_web) at (-6.0, 1.6) {
        \textbf{Bộ điều hợp Web / SSE}\\[2pt]
        Giao diện học giả, Luồng sự kiện
    };
    \node[adaptbox_dr, text width=3.6cm, minimum width=3.8cm, below=28pt of ad_web] (ad_api) {
        \textbf{Bộ điều hợp Tiếp nhận API}\\[2pt]
        Nhận tài liệu PDF \& Truy vấn
    };

    % --------------------------------------------------------------------------
    % Khối 2: Tầng Lõi Nghiệp Vụ & Cổng Giao Tiếp (Core & Ports) - Ở giữa
    % --------------------------------------------------------------------------
    \node[card_core, minimum width=5.6cm, minimum height=8.8cm] (c_core) at (0, -0.1) {};
    \node[align=center, font=\linespread{1.0}\selectfont\bfseries\scriptsize, text=blue!90!black] at (0, 3.8) {LÕI NGHIỆP VỤ \& CỔNG GIAO TIẾP\\[2pt](CORE DOMAIN \& PORTS)};
    
    % Cổng Phía Nhập (Inbound Port)
    \node[portbox, text width=4.6cm, minimum width=4.8cm] (p_in) at (0, 2.35) {
        \textbf{Cổng Phía Nhập (Inbound Port)}\\[2pt]
        Giao diện Phân tích \& Truy vấn Học thuật
    };

    % Lõi miền nghiệp vụ thuần khiết
    \node[compbox, text width=4.6cm, minimum width=4.8cm, fill=blue!8, below=15pt of p_in] (core) {
        \textbf{LÕI MIỀN NGHIỆP VỤ THUẦN KHIẾT}\\[2pt]
        Thực thể học thuật (\textit{Paper, Chunk, BBox}) $\parallel$ RRF
    };

    % Cổng Phía Xuất (Outbound Ports)
    \node[portbox, text width=4.6cm, minimum width=4.8cm, below=15pt of core] (p_out) {
        \textbf{Cổng Phía Xuất (Outbound Ports)}\\[2pt]
        • Cổng Trí tuệ AI \& Hộp cát MicroVM\\[1pt]
        • Cổng Lưu trữ Vector HNSW, BM25 \& CAS
    };

    % --------------------------------------------------------------------------
    % Khối 3: Tầng Điều Hợp Hạ Tầng (Driven Adapters) - Bên phải
    % --------------------------------------------------------------------------
    \node[card_driven, minimum width=4.5cm, minimum height=8.6cm] (c_driven) at (6.0, 0) {};
    \node[align=center, font=\linespread{1.0}\selectfont\bfseries\scriptsize, text=teal!90!black] at (6.0, 3.8) {BỘ ĐIỀU HỢP HẠ TẦNG\\[2pt](DRIVEN ADAPTERS)};
    
    \node[adaptbox_dn, text width=3.6cm, minimum width=3.8cm] (ad_llm) at (6.0, 2.0) {
        \textbf{Bộ điều hợp Dịch vụ LLM}\\[2pt]
        Mô hình ngôn ngữ lớn, Embedding
    };
    \node[adaptbox_dn, text width=3.6cm, minimum width=3.8cm, below=16pt of ad_llm] (ad_e2b) {
        \textbf{Bộ điều hợp Hộp cát E2B}\\[2pt]
        MicroVM Firecracker cô lập
    };
    \node[adaptbox_dn, text width=3.6cm, minimum width=3.8cm, below=16pt of ad_e2b] (ad_db) {
        \textbf{Bộ điều hợp CSDL \& CAS}\\[2pt]
        PostgreSQL 16, pgvector, CAS
    };

    % --------------------------------------------------------------------------
    % Đường kết nối & Luồng phụ thuộc hướng tâm
    % --------------------------------------------------------------------------
    % Driving Adapters -> Inbound Port
    \draw[arrowstyle] (ad_web.east) -- ++(0.4, 0) |- ([yshift=6pt]p_in.west);
    \draw[arrowstyle] (ad_api.east) -- ++(0.4, 0) |- ([yshift=-6pt]p_in.west);

    % Inbound Port -> Core
    \draw[arrowstyle] (p_in.south) -- (core.north);

    % Core -> Outbound Ports
    \draw[arrowstyle] (core.south) -- (p_out.north);

    % Driven Adapters -> Outbound Ports (Đảo ngược phụ thuộc DIP)
    \draw[dasharrow] (ad_llm.west) -- ++(-0.4, 0) |- ([yshift=14pt]p_out.east);
    \draw[dasharrow] (ad_e2b.west) -- ++(-0.4, 0) |- ([yshift=0pt]p_out.east);
    \draw[dasharrow] (ad_db.west) -- ++(-0.4, 0) |- ([yshift=-14pt]p_out.east);
\end{tikzpicture}
\caption{Mô hình phân rã module và cấu trúc Cổng -- Điều hợp (Ports and Adapters)}
\label{fig:hexagonal_arch}
\end{figure}

Luồng phụ thuộc trong hệ thống luôn có chiều hướng tâm: các bộ điều hợp kỹ thuật bên ngoài phụ thuộc vào các Cổng trong lõi, trong khi logic nghiệp vụ học thuật ở trung tâm hoàn toàn không biết gì về cơ sở hạ tầng thực tế bên dưới. Điều này bảo đảm tính liêm chính và sự ổn định cao độ cho toàn bộ nền tảng.
